\section{Research Questions}
\label{sec:3}

\subsection{The gap from sensing to decisions}

The cited literature establishes three foundations: multimodal XR assistance for individuals, physiological and behavioral analysis of collaborators, and mechanisms for moderation or scaffolding ~\cite{pro2,guide,p52,p57,p34,p51}. What remains to be established for the proposed setting is a complete comparison in which multi-person physiology changes assistance decisions, the comparison agent has the same behavioral and task information, and the resulting effects are measured at the appropriate team and individual levels. This formulation deliberately leaves room for a negative result: physiology may be redundant once task and behavioral context are available.

The key mechanism to test is relational rather than merely individual. A combined model might improve because it observes the target user's arousal, because it observes teammates' responses, or because it captures relationships among their signals. Existing interpersonal sensing and heterogeneity findings motivate distinguishing these explanations through ablation rather than attributing any gain to collaboration awareness ~\cite{p52,p32,hhiss}. Separate controls will test sensing value, relational value, and intervention timing.

\subsection{Research questions}

RQ1 Sensing value. Does adding physiological data improve prospective identification of a task-grounded team support need over speech, gaze, interaction traces, and task context alone?

RQ2 Relational value. Does information from teammates and their temporal relationships provide value beyond the target user's physiology? Where do disagreement, missingness, and shared task events explain errors?

RQ3 Intervention value. With the same available actions and comparable wording, does the physiology-informed policy improve appropriate information sharing or communication repair while limiting unwanted interruption?

RQ4 Transfer and context. If the core mechanism is validated, which parts remain useful in medical training and collaborative education, and do improvements persist in new tasks without agent prompts?

These questions are proposed tests. Their answers, effect directions, and attainable performance remain open.
